Networked Control Systems (NCS) demand low-latency, deterministic, and fault-tolerant communication, yet existing solutions such as Switched Ethernet suffer from single points of failure and lack multipath routing support. We introduce TinyNet, an open-source networking protocol for Networked Control Systems, implementable on system endpoints with minimal complexity. Operating without global state, TinyNet achieves memory efficiency and fault tolerance. It facilitates real-time multipath routing via a busyness-metric cost function that weights hop count against downstream buffer occupancy, reducing routing bottlenecks and distributing load across available paths. TinyNet targets small packet sizes to reduce transmission overhead, making it well-suited for control-loop traffic dominated by single-packet messages. \textcolor{red}{In simulations of a 16-node mesh, TinyNet sustains about 1.9 times the load of single-path shortest-path routing that is not planned for the traffic, matching the capacity of routes planned in advance. After one to four simultaneous node failures, its median blackout is 1.2--3.9~ms with at most 13 packets lost per flow. With two or more failures, loop-free-alternate fast reroute with the fastest documented failure detection typically loses hundreds of packets while awaiting routing reconvergence.} These results make TinyNet \textcolor{red}{a candidate} for distributed robotic control systems.